%!TEX root = ../thesis.tex

\section{Copied Bayesian notes}
\label{sec:change-point-copied-notes}

\noindent\textit{Draft transcription note.}
The source PDFs are handwritten image-only notebook exports. The page images below are included as facsimiles so that every page is preserved in the thesis draft. The text following each page is a first-pass transcription of the content that was legible during copying; uncertain words and symbols are marked explicitly.

\providecommand{\notepage}[4]{%
\clearpage
\subsection*{#1, page #2}
\begin{center}
\includegraphics[height=0.58\textheight,width=\linewidth,keepaspectratio]{chapter_change_point_detection/#3}
\end{center}
#4
}

\subsection{Source 1: \texttt{all-the-bayesian-math.pdf}}

\notepage{All-the-Bayesian-math}{1}{figures/note_pages/all_the_bayesian_math/page-01.png}{%
\begin{itemize}
\item Heading: ``Spontaneous meeting with Henning'', dated 04.06.19.
\item Questions: generation in direct and indirect pathways; how many animals does not work; capacity problem or time-scale issue.
\item Idea: does prefrontal capacity influence learning of semantic generalization?
\item Capacity is not comparable across different sizes of animals/persons.
\item Need to get generalization, i.e. generalization drops at a certain level.
\item Find point where generalization breaks. Fewer manipulable prefrontal capacity without changing representation, e.g. hidden layer.
\end{itemize}
}

\notepage{All-the-Bayesian-math}{2}{figures/note_pages/all_the_bayesian_math/page-02.png}{%
\begin{itemize}
\item Network 2: just best, certainly problematic.
\item Causality to semantic makes sense. What happens when episodic contradicts with semantic?
\item How are wrong mistakes in memory?
\item Retropale amnesia depends on lesion size.
\end{itemize}
}

\notepage{All-the-Bayesian-math}{3}{figures/note_pages/all_the_bayesian_math/page-03.png}{%
\begin{itemize}
\item Heading: ``Meeting Henning (Lund)'', dated 15.11.19.
\item Difference between reinforcement learning and association learning: value learning is stimulus-reward association learning.
\item Schema learning in deep nets: hidden-layer size works for perceptual learning; dynamics in different layers determine what other layers have already learned.
\item Hypothesis: first layer learns at one point, lower layers after, depending on what layer learned, such that the next layer can learn better after.
\item Generality transfer learning: use another net, starting point, and add a new layer that learns based on that schema.
\item Model as meta-learning, autoencoder with bottleneck, new episodes learned only in hippocampus; the rest is trained to make learning in hippocampus easier.
\item Semantic development in children: representations in hippocampus are better suited for learning associations; input pathway develops further.
\end{itemize}
}

\notepage{All-the-Bayesian-math}{4}{figures/note_pages/all_the_bayesian_math/page-04.png}{%
\begin{itemize}
\item MRI study with Yee Lee: test statistical learning with paired associations and different category-correlation strengths.
\item Task sketch: shape parameter/cue 1 (houses), response 1, circles, shape parameter/cue 2 (trees), response 2.
\item Hierarchical structure: hippocampal layer $\eta_1>\eta_2>\eta_3>\eta_4>\eta_5$.
\item Learning-specific cue-response associations and global house-to-squares association, probably parameters $P_1$ and $P_2$.
\item Hypothesis: given weak connection $p$, all layers with $\eta<\eta^\ast$ can learn statistics well. For different $p$, different layers should be involved in the recall process; perhaps all layers are involved but some just do not produce useful answers.
\end{itemize}
}

\notepage{All-the-Bayesian-math}{5}{figures/note_pages/all_the_bayesian_math/page-05.png}{%
\begin{itemize}
\item Response paths are drawn from a Bernoulli distribution. Compute posterior distribution $q$.
\item Samples enter the model with connection probability $p$. $T_i$ are number of samples layer $i$ can store and use for estimate $\hat p_i$.
\item If layers learn at same speed but remember only $T_i$ samples, the smaller $T_i$ gives worse estimate.
\item If $\eta_i>\eta_{i+1}$ and $T_i<T_{i+1}$, deeper layers learn faster and also remember less or forget faster. The best estimate of a layer should also be reached slower for deeper layers.
\item Proposed readout:
\[
q(P_i,t\mid \eta_i,T_i)
\]
and question whether the brain uses $\hat q(\hat P_i,\hat\eta_i,t,\hat T_i)$.
\end{itemize}
}

\notepage{All-the-Bayesian-math}{6}{figures/note_pages/all_the_bayesian_math/page-06.png}{%
\begin{itemize}
\item Heading: ``Calculating the posterior (best estimate of $p$ that a given shortcut could learn from its observations)''.
\item Random variable $R\in\{0,1\}$ is a response.
\item Likelihood of single trial:
\[
\Pr(R=1)=p,\qquad \Pr(R;p)=
\begin{cases}
p,&R=1,\\
1-p,&R=0,
\end{cases}
=p^R(1-p)^{R-1}.
\]
\item Likelihood after $n$ binary cues:
\[
X\sim \operatorname{Bin}(n,p),\qquad
\Pr(X=k;n,p)=\binom{n}{k}p^k(1-p)^{n-k}.
\]
\item Conjugate prior for Bernoulli/binomial:
\[
\Pr(p;\alpha,\beta)=Zp^{\alpha-1}(1-p)^{\beta-1},
\quad
Z=\frac{1}{B(\alpha,\beta)}
=\frac{\Gamma(\alpha+\beta)}{\Gamma(\alpha)+\Gamma(\beta)}.
\]
\item Posterior:
\[
\Pr(p=x\mid k)\propto x^{\alpha_0+k-1}(1-x)^{n-k+\beta_0-1}
=\operatorname{Beta}(\alpha_0+k,\beta_0+n-k).
\]
\end{itemize}
}

\notepage{All-the-Bayesian-math}{7}{figures/note_pages/all_the_bayesian_math/page-07.png}{%
\begin{itemize}
\item Posterior rewritten:
\[
\Pr(p;k,n)=
\frac{x^{k+\alpha_0-1}(1-p)^{n-k+\beta_0-1}}
{B(k+\alpha_0,n-k+\beta_0)}
=\operatorname{Beta}(\alpha_0+k,\beta_0+n-k).
\]
\item How to get $\eta_i$ and $T_i$ into $Z$:
\begin{enumerate}
\item First ignore learning rate $\eta_i$; shortcuts just have different capacity $C_i$, so $k_i$ are the $R=1$ responses in the last $n_i$ observations.
\item Total time $t$ is all shown samples. Compile $\operatorname{Beta}(\alpha',\beta')$ with
\[
\alpha'=\alpha_0+\sum_{i=1}^{t} F(i)\delta(R_i-1),
\qquad
\beta'=\beta_0+\sum_{i=1}^{t} F(i)\delta(R_i),
\]
where $F(i)=\Theta(i-(t-C_j))$ or $F(i)=e^{-\tau_j(t-i)}$.
\end{enumerate}
\end{itemize}
}

\notepage{All-the-Bayesian-math}{8}{figures/note_pages/all_the_bayesian_math/page-08.png}{%
\begin{itemize}
\item Heading: ``Multiple response classes ($>2$)''.
\item Likelihood becomes categorical distribution and conjugate prior becomes Dirichlet distribution.
\item If $N_r$ is number of response classes:
\[
\operatorname{Dir}(\vec\alpha)=\frac{1}{B(\vec\alpha)}
\prod_{r=1}^{N_r}x_r^{\alpha_r-1},
\qquad
\vec\alpha=(\alpha_1,\ldots,\alpha_{N_r}).
\]
\item Position/posterior:
\[
\alpha'_r=\alpha^0_r+\sum_{i=1}^{T}F(i)\delta(R(i)-r),
\]
with $F(t)=\Theta(t-(T-C_j))$ or $F(t)=e^{-\tau_j(T-t)}$.
\item Multiple cues: now $t$ is current observation index; $\eta_q(t)$ is number of observations with cue $q$ until $t$ and $\sum_q \eta_q(T)=T$. $R_t,Q_t$ are response and cue at observation $t$.
\end{itemize}
}

\notepage{All-the-Bayesian-math}{9}{figures/note_pages/all_the_bayesian_math/page-09.png}{%
\begin{itemize}
\item Heading: ``Multiple cues with marginal posterior [Beta$(\alpha_r,\sum_i\alpha_i-\alpha_r)$]''.
\item $t$ is current observation index. $\eta_q(t)$ is number of observations with cue $q$ until $t$, $\sum_q\eta_q(T)=T$.
\item $R_t,Q_t$ are response and cue at observation $t$.
\item Posterior per cue:
\[
\operatorname{Beta}(\alpha^q_1,\beta^q_1),\qquad
\alpha^q_r=\alpha^q_0+\sum_{i=1}^{t}F(i)\delta(R_i-r)\delta(Q_i-q).
\]
\item Equivalently,
\[
\beta^q_r=\sum_{r=1}^{N_r}\alpha^q_r-\alpha^q_r
=(N_r-1)\alpha_0+\sum_{r=1}^{N_r-1}\sum_{i=1}^{t}F(i)\delta(Q_i-q)\delta(R_i-r).
\]
\item Red annotation: associated response is $r=0$.
\end{itemize}
}

\notepage{All-the-Bayesian-math}{10}{figures/note_pages/all_the_bayesian_math/page-10.png}{%
\begin{itemize}
\item Heading: ``Readout: combining shortcut posteriors'', dated 3.2.20.
\item Let $N_s$ shortcuts each compute a different posterior over $p$:
\[
P_s^t=P_s(p\mid\{R_i\}_{i=1}^t)=\operatorname{Beta}(\alpha',\beta'),
\]
\[
\alpha'=\alpha_0+\sum_{i=1}^{t}F_s(i)\delta(R_i-1),\qquad
\beta'=\beta_0+\sum_{i=1}^{t}F_s(i)\delta(R_i).
\]
\item How would an optimal readout combine posteriors?
\begin{enumerate}
\item Weighted sum:
\[
P_\Sigma^t=\sum_{s=1}^{N_s}C_s(t)P_s^t,
\]
with $C_s(t)$ chosen by comparing with true posterior.
\item Product:
\[
P_\Pi^t=\prod_{s=1}^{N_s}P_s^t.
\]
\end{enumerate}
\item Issue: assumes statistical independence, but posteriors depend on the same observations.
\end{itemize}
}

\notepage{All-the-Bayesian-math}{11}{figures/note_pages/all_the_bayesian_math/page-11.png}{%
\begin{itemize}
\item Heading: ``Random switching of association $P$''.
\item $P$ switches approximately with Poisson switching; inter-switch interval is exponential.
\item Transition probability:
\[
P_{\mathrm{trans}}(t)=P(P_t\mid P_{t-1})
=\alpha\,U(0,1)+(1-\alpha)\delta(P_{t-1}-P_t).
\]
\item Prior at $t=1$: $P(P_1)=U(0,1)=\operatorname{Beta}(1,1)$.
\item Example posterior update:
\[
P(R_2=1\mid P_2,P_1)
=\alpha\,\operatorname{Beta}(\alpha_2^1,\beta_2^1)
(1-\alpha)\operatorname{Bern}(P_2)\delta(P_2-P_1).
\]
\item Bayes theorem box:
\[
P(p\mid R_i)=\frac{P(R_i\mid p)P(p)}{P(R_i)}.
\]
\end{itemize}
}

\notepage{All-the-Bayesian-math}{12}{figures/note_pages/all_the_bayesian_math/page-12.png}{%
\begin{itemize}
\item Page contains a tree/table expansion for full posterior, not normalized.
\item First level:
\[
P(p_1)=U(0,1)=\operatorname{Beta}(1,1),\quad
P(R_1\mid p_1)=\operatorname{Bern}(p_1),\quad
P(p_1\mid R_1)=\operatorname{Beta}(\alpha_1^1,\beta_1^1).
\]
\item State counters:
\[
\alpha_i^j=1+\sum_{k=i}^{j}\delta(R_k-1),\qquad
\beta_i^j=1+\sum_{k=i}^{j}\delta(R_k).
\]
\item Full posterior sketched as a sum over switch histories:
\[
P(p_n\mid\{R_i\}_{i=1}^{n})
=\sum_{m=1}^{n}\alpha^{\,n-m}(1-\alpha)^{m-1}
\cdots
\]
with products of Beta terms and delta constraints linking equal $p$ values between switches.
\end{itemize}
}

\notepage{All-the-Bayesian-math}{13}{figures/note_pages/all_the_bayesian_math/page-13.png}{%
\begin{itemize}
\item Heading: ``Normalize posteriors''.
\item Beta density and normalizer:
\[
\operatorname{Beta}(\alpha,\beta)=
\frac{p^{\alpha-1}(1-p)^{\beta-1}}{B(\alpha,\beta)},\qquad
B(\alpha,\beta)=\int_0^1 q^{\alpha-1}(1-q)^{\beta-1}\,dq.
\]
\item Product identity:
\[
\operatorname{Beta}(\alpha,\beta)\operatorname{Beta}(\gamma,\delta)
\propto \operatorname{Beta}(\alpha+\gamma-1,\beta+\delta-1).
\]
\item Boxed special case:
\[
\operatorname{Beta}(\alpha_i^j,\beta_i^j)\operatorname{Beta}(\alpha_k^\ell,\beta_k^\ell)
=\operatorname{Beta}(\alpha_i^j+\alpha_k^\ell-1,\beta_i^j+\beta_k^\ell-1)
\]
for matching intervals/indices.
\end{itemize}
}

\notepage{All-the-Bayesian-math}{14}{figures/note_pages/all_the_bayesian_math/page-14.png}{%
\begin{itemize}
\item Heading: ``Normalize posteriors contd.''
\item Integral of Beta times deltas:
\[
\int_0^1 \operatorname{Beta}(\alpha_i^j,\beta_i^j)\delta(q_j-p_i)\,dq_j
\]
leading to terms $p_i^{\alpha_i^j-1}(1-p_i)^{\beta_i^j-1}$ and cancellation of $B(\alpha_i^j,\ldots)$.
\item Counting observations:
\[
\alpha_i^j=1+\sum_{m=i}^{j}\delta(R_m-1),
\]
and for total observations $n=j-i+1$ and $k=\sum_{m=i}^{j}\delta(R_m-1)$,
\[
2p_i^k(1-p_i)^{n-k}
=\operatorname{Binom}(n,p_i\mid k)\frac{1}{\binom{n}{k}}.
\]
\item Product of Bernoulli terms:
\[
\operatorname{Bn}(p\mid k)\operatorname{Bn}(p\mid \ell)
=p^{k+\ell}(1-p)^{2-(k+\ell)}
=\operatorname{Binom}(2,p\mid k+\ell)\frac{1}{\binom{2}{k+\ell}},
\quad k,\ell\in\{0,1\}.
\]
\end{itemize}
}

\notepage{All-the-Bayesian-math}{15}{figures/note_pages/all_the_bayesian_math/page-15.png}{%
\begin{itemize}
\item Page repeats the switching prior tree with $\operatorname{Beta}_i^j(p)=\operatorname{Beta}(p,\alpha_i^j,\beta_i^j)$.
\item For $n=2$, one branch gives:
\[
P(p_2\mid R_2,R_1)=
\operatorname{Beta}_1^2(p_2)\left[\alpha+(1-\alpha)\delta(p_2-p_1)\right],
\]
with normalization via
\[
\int_0^1 \operatorname{Beta}(\alpha_i^j,\ldots)\delta(q_j-p_i)\,dq_j
=p_i^k(1-p_i)^{n-k}
=\frac{1}{\binom{n}{k}}\operatorname{Binom}(n,p_i\mid k).
\]
\end{itemize}
}

\notepage{All-the-Bayesian-math}{16}{figures/note_pages/all_the_bayesian_math/page-16.png}{%
\begin{itemize}
\item Heading: ``If switch, use flat prior.''
\item Prior/likelihood/posterior tree expanded for $n=1,2,3,4$ with switch probability $\alpha$ and no-switch probability $1-\alpha$.
\item Intermediate terms for $P(p_3\mid R\ldots)$ are expanded into powers of $\alpha$ and $(1-\alpha)$.
\item Red-boxed generalization:
\[
P(p_n\mid R\ldots)=\beta_1(1-\alpha)^{n-1}
\sum_{i=2}^{n}\beta_i
\sum_{m=n-i+1}^{n-1}\binom{i-2}{t-m-1}\alpha^{n-m}(1-\alpha)^{m-1}
\]
where the exact handwritten indices are ambiguous.
\end{itemize}
}

\notepage{All-the-Bayesian-math}{17}{figures/note_pages/all_the_bayesian_math/page-17.png}{%
\begin{itemize}
\item Factor derivation:
\[
P(p_t\mid\{R_i\}_{i=0}^{t})
=\beta_{1,t}(1-\alpha)^{t-1}
\sum_{i=2}^{t}\beta_{i,t}\sum_{m=t-i+1}^{t-1}
\binom{i-2}{t-m-1}\alpha^{t-m}(1-\alpha)^{m-1}.
\]
\item Listed factors:
\[
\beta_{1,t}:(1-\alpha)^{t-1},
\]
\[
\beta_{2,t}:\sum_{m=t-1}^{t-1}\binom{0}{t-1-m}\alpha^{t-m}(1-\alpha)^{m-1}
=\alpha(1-\alpha)^{t-2},
\]
and analogous expressions for $\beta_{3,t}$ and $\beta_{4,t}$.
\item Boxed general form:
\[
P(p_t\mid\{R_i\}_{i=0}^{t})
=\beta_{0,t}(1-\alpha)^t
\sum_{i=1}^{t}\beta_{i,t}\sum_{m=1}^{i}
\binom{i-1}{m-1}\alpha^m(1-\alpha)^{(t-1)-(\cdots)}.
\]
\end{itemize}
}

\notepage{All-the-Bayesian-math}{18}{figures/note_pages/all_the_bayesian_math/page-18.png}{%
\begin{itemize}
\item Posterior represented as mixture:
\[
P(p_t\mid\{R_i\}_{i=0}^{t})=\sum_{i=0}^{t} f_i(t)\beta_i(t).
\]
\item Recursive/closed form weights:
\[
f_i(t)=
\begin{cases}
(1-\alpha)^t,&i=0,\\
\sum_{i=1}^{t}\beta_i(t)\sum_{m=1}^{i}\binom{i-1}{m-1}\alpha^m(1-\alpha)^{t-m},&\text{else}.
\end{cases}
\]
\item Examples:
\[
f_1(t)=\alpha(1-\alpha)^{t-1},\quad
f_2(t)=\alpha(1-\alpha)^{t-1}+\alpha^2(1-\alpha)^{t-2},
\]
\[
f_t(t)=\sum_{m=1}^{t}\binom{t-1}{m-1}\alpha^m(1-\alpha)^{t-m}=\alpha.
\]
\item Initial term $f_0(t)=1-\alpha$ is noted in the lower red box.
\end{itemize}
}

\notepage{All-the-Bayesian-math}{19}{figures/note_pages/all_the_bayesian_math/page-19.png}{%
\begin{itemize}
\item Markovian likelihood note. For $t=3$:
\[
P(p_3)=P(p_2\mid p_1)\{R_i\}_{i=1}^{2},
\]
\[
P(R_3=1\mid p_3,p_2)=P(R_3=1\mid p_3)P(p_3\mid p_2)
\]
and terms with switch/no-switch:
\[
\alpha\operatorname{Beta}(\alpha_1',\beta_1')
(1-\alpha)\delta(p_3-p_2)\operatorname{Bern}(p_3).
\]
\item Posterior expanded as products of the previous posterior and current likelihood. Red note says ``TODO: check!''.
\item Result shows terms in $\alpha^2$, $\alpha(1-\alpha)$, and $(1-\alpha)^2$ with products of Beta densities and delta functions.
\end{itemize}
}

\notepage{All-the-Bayesian-math}{20}{figures/note_pages/all_the_bayesian_math/page-20.png}{%
\begin{itemize}
\item Generalized version of page 19 for $t=n$:
\[
P(p_n)=P(p_n\mid p_{n-1})\{R_i\}_{i=1}^{n-2},
\qquad
P(R_n=1\mid p_n,p_{n-1})=P(R_n=1\mid p_n)P(p_n\mid p_{n-1}).
\]
\item Transition likelihood:
\[
P(R_n=1\mid p_n)P(p_n\mid p_{n-1})
=\alpha\operatorname{Beta}(\alpha_1',\beta_1')
(1-\alpha)\delta(p_n-p_{n-1})\operatorname{Bern}(p_n).
\]
\item Posterior:
\[
P(p_n\mid p_{n-1},R_n)=
\alpha^{n-1}[\operatorname{Beta}(\alpha_1',\beta_1')]^n
(1-\alpha)\left[\prod_{i=1}^{n-1}\delta(p_{i+1}-p_i)\right]\operatorname{Beta}(\alpha_n',\beta_n')
\cdots
\]
\item Lower part copies from $t=3$ and simplifies terms with $\delta(p_3-p_2)$ and $\delta(p_2-p_1)$.
\end{itemize}
}

\notepage{All-the-Bayesian-math}{21}{figures/note_pages/all_the_bayesian_math/page-21.png}{%
\begin{itemize}
\item Normalization constants $Z_i^j$ for products of Bernoulli and Beta terms.
\item Prior $p=\beta_0$, likelihood $\ell=\operatorname{Bn}(p)$ and posterior $\beta_2^2=p\ell/Z$.
\item Identities:
\[
Z_2^2=\frac{B(\alpha_2^2,\beta_2^2)}{B(1,1)},\qquad
Z_1^2=\frac{B(\alpha_1^2,\beta_1^2)}{B(\alpha_1^1,\beta_1^1)}.
\]
\end{itemize}
}

\notepage{All-the-Bayesian-math}{22}{figures/note_pages/all_the_bayesian_math/page-22.png}{%
\begin{itemize}
\item Continues normalization:
\[
P(p_2\mid\{R_i\})=
\left[\alpha Z_2^2\beta_2^2(p)+(1-\alpha)Z_1^2\beta_1^2(p)\right]\frac{1}{Z}.
\]
\item Normalizing integral:
\[
Z=\int_0^1 dp\left[
\alpha\frac{B(\alpha_2^2,\beta_2^2)}{B(1,1)}\beta_2^2(p)
+(1-\alpha)\frac{B(\alpha_1^2,\beta_1^2)}{B(\alpha_1^1,\beta_1^1)}\beta_1^2(p)
\right].
\]
\item Because each beta density integrates to one:
\[
P(p_2\mid\{R_i\})=
\frac{\alpha Z_2^2\beta_2^2(p)+(1-\alpha)Z_1^2\beta_1^2(p)}
{\alpha Z_2^2+(1-\alpha)Z_1^2}.
\]
\item For $t=3$, analogous mixture weights are derived, ending in a rational expression with $\alpha Z_3^3$, $\alpha(1-\alpha)Z_2^3Z_2^2$, and $(1-\alpha)^2Z_1^3Z_1^2$ terms.
\end{itemize}
}

\notepage{All-the-Bayesian-math}{23}{figures/note_pages/all_the_bayesian_math/page-23.png}{%
\begin{itemize}
\item Heading: ``Bernoulli--Beta conjugate pair''.
\item Bernoulli:
\[
X\sim\operatorname{Bernoulli}(p),\qquad P(X=1)=p,
\]
\[
\operatorname{Bn}(k\mid p)=\operatorname{Bn}(X=k\mid p)=p^k(1-p)^{1-k},\quad k\in\{0,1\}.
\]
\item Beta:
\[
X\in[0,1],\quad
\operatorname{Beta}(x\mid\alpha,\beta)
=x^{\alpha-1}(1-x)^{\beta-1}
\left[\int_0^1 u^{\alpha-1}(1-u)^{\beta-1}du\right]^{-1}
=\frac{x^{\alpha-1}(1-x)^{\beta-1}}{B(\alpha,\beta)}.
\]
\item Posterior from Bernoulli likelihood and Beta prior:
\[
P(p\mid k)=\frac{P(k\mid p)P(p)}{\int_0^1P(k\mid q)P(q)dq}
=\operatorname{Beta}(k+\alpha_0,1-k+\beta_0).
\]
\end{itemize}
}

\notepage{All-the-Bayesian-math}{24}{figures/note_pages/all_the_bayesian_math/page-24.png}{%
\begin{itemize}
\item Heading: ``Henning notes / Underlying graphical model''.
\item Variables: time since last change in contingency $n_t$, association strength $s_t$, observations $o_t$. Graphical sketch: $n_t\to n_{t+1}$, $n_t\to s_{t+1}$, $s_t\to s_{t+1}$, $s_t\to o_t$.
\item Object of interest:
\[
P(s,n;t\mid o_{1:t})
\]
as joint probability that last change of contingency happened $n$ time steps ago and association strength is $s$.
\item Marginal:
\[
P(s;t\mid o_{1:t})=\sum_n P(s,n;t\mid o_{1:t}).
\]
\item Update rule for posterior:
\[
P(s,n;t+1\mid o_{1:t+1})
=\frac{1}{Z}P(o_{t+1}\mid s_{t+1})P(s,n;t\mid o_{1:t})
\]
with an integral/sum over prior states and transition probability $P(s,n\mid s',n')$.
\end{itemize}
}

\notepage{All-the-Bayesian-math}{25}{figures/note_pages/all_the_bayesian_math/page-25.png}{%
\begin{itemize}
\item Heading: ``Choices for $P(s,n\mid s',n')$''.
\item Perfect integrator:
\[
P(s,n\mid s',n')=\delta_{n,n'+1}\delta(s-s').
\]
World never changes.
\item Leaky/exponential model:
\[
P(s,n\mid s',n')=(1-\alpha)\delta_{n,n'+1}\delta(s-s')+\alpha\delta_{n,0}U(s).
\]
Probability that world changes is constant for all time steps. Switch probability can depend on how long ago the last switch occurred; allows long survival times, e.g. power law for $\alpha(n)=k/n$.
\item Hippocampus as a cascade model: hidden states $n=0,1,2,3$ with transition probabilities $1-\gamma^{(0)},1-\gamma^{(1)},1-\gamma^{(2)}$ and links to cortex through $\alpha^{(1)},\alpha^{(2)},\alpha^{(3)}$.
\end{itemize}
}

\notepage{All-the-Bayesian-math}{26}{figures/note_pages/all_the_bayesian_math/page-26.png}{%
\begin{itemize}
\item $t=1$ initialization:
\[
P(s)=U(s)=\operatorname{Beta}(1,1)=\beta_0,\quad s\in[0,1],
\]
\[
P(o\mid s)=\operatorname{Bn}(s),\qquad
P(s\mid o)=\operatorname{Bn}(s)\beta_0=\beta_1^{(1)}.
\]
\item At time $t$:
\[
P(s\mid o_{1:t})=\sum_{i=1}^{t}f_i(t,n,o_{1:t})\beta_i(t,s,o_{1:t}).
\]
\item Beta component:
\[
\beta_i(t)=\operatorname{Beta}(a_i(t),b_i(t)),
\]
\[
a_i(t)=1+\sum_{j=0}^{i-1}\delta(o_{t-j}-1),\quad
b_i(t)=1+\sum_{j=0}^{i-1}\delta(o_{t-j}).
\]
\item Update:
\[
\operatorname{Bn}(s,o_{t+1})\beta_i(t,s,o_{1:t})
=\beta_{i+1}(t+1,s,o_{1:t+1})Z_{i+1}(t+1),
\]
\[
Z_{i+1}(t+1)=\frac{B(a_{i+1}(t+1),b_{i+1}(t+1))}
{B(a_i(t),b_i(t))}.
\]
\end{itemize}
}

\notepage{All-the-Bayesian-math}{27}{figures/note_pages/all_the_bayesian_math/page-27.png}{%
\begin{itemize}
\item Continuation of update rule:
\[
P(s,n;t+1\mid o_{1:t+1})
=\frac{1}{Z}P(o_{t+1}\mid s_{t+1})P(s,n;t\mid o_{1:t})
\]
with transition probability and old posterior.
\item General transition:
\[
P(s,n\mid s',n')
=(1-\alpha(n'))\delta_{n,n'+1}\delta(s-s')+\alpha(n')\delta_{n,0}U(s).
\]
\item Likelihood:
\[
P(o_t\mid s_t)=\operatorname{Bn}(o_t;s).
\]
\item Prior/old posterior:
\[
P(s,n;t\mid o_{1:t})
=\sum_{i=1}^{t}f_i(t,n,o_{1:t})\beta_i(t,s,o_{1:t}).
\]
\item Resulting posterior expression separates no-change terms and new-change terms with $\alpha(n')$ and $(1-\alpha(n'))$.
\end{itemize}
}

\notepage{All-the-Bayesian-math}{28}{figures/note_pages/all_the_bayesian_math/page-28.png}{%
\begin{itemize}
\item Boxed expression:
\[
P(s,n;t\mid o_{1:t})
=\frac{1}{Z}\sum_{i=1}^{t}f_i(t-j,n,o_{1:t})\beta_i(t,s,o_{1:t}).
\]
\item For $i\ge2$:
\[
f_i(t,n,o_{1:t})=(1-\alpha(n-1))Z_i(t)f_{i-1}(t-1,n-1,o_{1:t-1})
\]
with
\[
Z_i(t)=\frac{B(a_i(t),b_i(t))}{B(a_{i-1}(t-1),b_{i-1}(t-1))}.
\]
\item New state normalization:
\[
Z=\int_0^1 ds'\,P(s,n;t+1\mid o_{1:t+1})
=\sum_{i=1}^{t}f_i(t,n,o_{1:t}).
\]
\item Initialization:
\[
P(s,n;t=1\mid o_1)=\frac{1}{Z}f_1(t=1,n=0,o_1)\beta_1(t=1,s,o_1).
\]
\end{itemize}
}

\notepage{All-the-Bayesian-math}{29}{figures/note_pages/all_the_bayesian_math/page-29.png}{%
\begin{itemize}
\item Repeats underlying graphical model, now with observation $x_t$ rather than $o_t$.
\item Object of interest:
\[
P(s_t,n_t\mid x_{1:t})=\sum_t P(s_t,n_t\mid x_{1:t})
\]
and posterior probability of association strength $s$.
\item Update rule:
\[
P(s_t,n_t\mid x_{1:t})
=\frac{1}{Z}P(x_t\mid s_t)P(s_t,n_t\mid x_{1:t-1})
\]
with transition probability and old posterior:
\[
\sum_{n_{t-1}}\int ds_{t-1}\,
P(s_t,n_t\mid s_{t-1},n_{t-1})P(s_{t-1},n_{t-1}\mid x_{1:t-1}).
\]
\item Notes: factorization from graphical model, conjugate priors/media.
\end{itemize}
}

\notepage{All-the-Bayesian-math}{30}{figures/note_pages/all_the_bayesian_math/page-30.png}{%
\begin{itemize}
\item Heading: ``Choices for $P(s_t,n_t\mid s_{t-1},n_{t-1})$''.
\item Perfect integrator:
\[
P(s_t,n_t\mid s_{t-1},n_{t-1})
=\delta_{n_t,n_{t-1}+1}\delta(s_t-s_{t-1}).
\]
\item Exponential/leaky model:
\[
P(s_t,n_t\mid s_{t-1},n_{t-1})
=(1-\gamma)\delta_{n_t,n_{t-1}+1}\delta(s_t-s_{t-1})
\gamma\delta_{n_t,0}U(s_t).
\]
\item More general model:
\[
P(s_t,n_t\mid s_{t-1},n_{t-1})
=(1-\gamma(n_{t-1}))\delta_{n_t,n_{t-1}+1}\delta(s_t-s_{t-1})
\gamma(n_{t-1})\delta_{n_t,0}U(s_t).
\]
\item Allows long survival times, e.g. power law $\gamma(n_t)=k/n_t$. Includes the same hippocampus-to-cortex cascade sketch as page 25.
\end{itemize}
}

\notepage{All-the-Bayesian-math}{31}{figures/note_pages/all_the_bayesian_math/page-31.png}{%
\begin{itemize}
\item Heading: ``Update rule for the posterior''.
\item New posterior:
\[
P(s_t,n_t\mid x_{1:t})
=\frac{1}{Z}P(x_t\mid s_t)P(s_t,n_t\mid x_{1:t-1}).
\]
\item Transition probability:
\[
P(s_t,n_t\mid s_{t-1},n_{t-1})
=(1-\gamma(n_{t-1}))\delta_{n_t,n_{t-1}+1}\delta(s_t-s_{t-1})
\gamma(n_{t-1})\delta_{n_t,0}U(s_t).
\]
\item Likelihood $P(x_t\mid s_t)=\operatorname{Bn}(x_t\mid s_t)$.
\item Prior/old posterior:
\[
P(s_{t-1},n_{t-1}\mid x_{1:t-1})
=\sum_{i=1}^{t-1}f_i(n_{t-1}\mid x_{1:t-1})\beta_i(s_{t-1}\mid x_{1:t-1}).
\]
\item Posterior contains a normalized expression with two sums: no-change terms using $(1-\gamma(n_{t-1}))$ and reset terms using $\gamma(n_{t-1})$.
\end{itemize}
}

\notepage{All-the-Bayesian-math}{32}{figures/note_pages/all_the_bayesian_math/page-32.png}{%
\begin{itemize}
\item Red-boxed unnormalized posterior:
\[
P'(s_t,n_t\mid x_{1:t})
=\sum_{i=1}^{t} f_i(n_t\mid x_{1:t})\beta_i(s_t\mid x_{1:t}).
\]
\item Reindexed no-change term:
\[
\left[\sum_{i=2}^{t}(1-\gamma(n_t-1))f_{i-1}(n_t-1\mid x_{1:t-1})
Z_i(x_t)y(u-1)\beta_i(s_t\mid x_{1:t})\right]
\]
plus reset term with $\gamma(n_{t-1})f_i(\cdots)$, $\delta_{n_t,0}$, $Z_0(x_{1:t})$, and $\beta_1(s_t\mid x_{1:t})$.
\item Normalization:
\[
Z=\sum_{n_t=0}^{t-1}\int_0^1 ds_t\,P'(s_t,n_t\mid x_{1:t})
=\sum_{n_t=0}^{t-1}\sum_{i=1}^{t} f_i(n_t\mid x_{1:t}).
\]
\item Marginalized posterior:
\[
P(s_t\mid x_{1:t})=\sum_{n_t=0}^{t}P(s_t,n_t\mid x_{1:t})
=\sum_{n_t=0}^{t}\sum_{i=1}^{t}f_i(n_t\mid x_{1:t})\beta_i(s_t\mid x_{1:t}).
\]
\item At bottom: for $t=1,n_t=0$, $P(s,n\mid x)=f_1(n=1\mid x_1)\beta_1(s_1\mid x_1)$.
\end{itemize}
}

\notepage{All-the-Bayesian-math}{33}{figures/note_pages/all_the_bayesian_math/page-33.png}{%
\begin{itemize}
\item Heading: ``In Einstein notation''.
\item Posterior written as:
\[
P(s_t,n_t\mid x_{1:t})
=\sum_{i=1}^{t} f_i(n_t\mid x_{1:t})\beta_i(s_t\mid x_{1:t}).
\]
\item Need tensors for $t_i,n_t,t_s\to t_{sn},f,\beta$; $f=\beta=\delta_{i,j}$ is noted.
\item Initial condition: $P(0;0)=1$.
\item Normalizer:
\[
Z=\sum_{n_t=0}^{t-1}\sum_{i=1}^{t}f_i(n_t\mid x_{1:t}),
\qquad
Z=\operatorname{einsum}(\text{``tin->t''},f)\quad\text{[in sum()]}.
\]
\item Recursion:
\[
f_1(n_t\mid x_{1:t})
=\frac{1}{Z}\left(\sum_{n_{t-1}=0}^{t-1}\gamma(n_{t-1})
\sum_{j=1}^{t-1}f_j(n_{t-1}\mid x_{1:t-1})Z_0\delta_{n_t,0}\right).
\]
\item Last line:
\[
f[t,1,0]=\frac{1}{2}\operatorname{einsum}(\text{``n,jn->''},\gamma(\operatorname{arange}(t)),f[t-1])\cdot Z_0\cdot \delta_{n_t,0}.
\]
\end{itemize}
}

\subsection{Source 2: \texttt{2025-01-meeting-francesco-AND-multinomial-math.pdf}}

\notepage{Francesco/multinomial notes}{1}{figures/note_pages/francesco_multinomial_math/page-1.png}{%
\begin{itemize}
\item Heading: ``Slides from Francesco'', dated 14.01.25.
\item Character $\to$ category association. Examples: A/B mapped to furniture, food, utensils, toys.
\item Question: what category given clue? Response: one item of the category, with category contingency roughly 75\%--25\%.
\item Aim: test episodic memory after interim stimulus. Before run also change probability?
\item Five blocks, one day? Three change points. About 24 trials per clue per block.
\end{itemize}
}

\notepage{Francesco/multinomial notes}{2}{figures/note_pages/francesco_multinomial_math/page-2.png}{%
\begin{itemize}
\item Questions for Francesco:
\begin{itemize}
\item What data exists? Any fMRI?
\item What is needed for a minimal paper?
\item Do you want to detect change points per participant based on their behaviour?
\end{itemize}
\item Need to look into papers again: Nassar, Wilson.
\item Variable per participant: likely fit for each participant.
\end{itemize}
}

\notepage{Francesco/multinomial notes}{3}{figures/note_pages/francesco_multinomial_math/page-3.png}{%
\begin{itemize}
\item Data structure: full behavior dataset, full train/test.
\item High-resolution ALC? subfields.
\item Agreements:
\begin{itemize}
\item Francesco sends behavior data, prototype, same category/item pair for each clue/participant.
\item fMRI: different pairs per participant; without MRI protocol; Nassar/Wilson; Dario/Nesa? current model from multiple-choice task (2-to-2) to this task (2-to-4 categories), needs Dirichlet instead of Bernoulli.
\item Later: fit different models, Nassar/Wilson; run their analyses in cascade.
\end{itemize}
\end{itemize}
}

\notepage{Francesco/multinomial notes}{4}{figures/note_pages/francesco_multinomial_math/page-4.png}{%
\begin{itemize}
\item Model to human behavior, fit per participant.
\item Not based on rule/truth; includes times, level learning rate.
\item Check papers for individual human differences (Nassar/Wilson).
\end{itemize}
}

\notepage{Francesco/multinomial notes}{5}{figures/note_pages/francesco_multinomial_math/page-5.png}{%
\begin{itemize}
\item Heading: ``Multinomial extension for $k=4$ answers'', dated 19.3.25.
\item Multinomial PMF:
\[
f(x_1,x_2,x_3,x_4;n,p_1,p_2,p_3,p_4)
=P(X_1=x_1,\ldots,X_4=x_4)
=\frac{n!}{\prod_i x_i!}\prod_i p_i^{x_i},
\]
where $\sum_i x_i=n$.
\item Extracting $n$ balls of different colors, replacing after each draw; $x_i$ is number of draws of ball color $i$.
\end{itemize}
}

\notepage{Francesco/multinomial notes}{6}{figures/note_pages/francesco_multinomial_math/page-6.png}{%
\begin{itemize}
\item Heading: ``Variable Bayes model''.
\item Parameters of multinomial with $k=4$:
\[
P_1,P_2,P_3,P_4,\qquad \sum_i P_i=1.
\]
\item Update rule for posterior:
\[
P(\vec P,s,n)
=\frac{1}{Z}P(X_t\mid \vec P,s)\,P(s)
\]
with intended dependence on latest changepoint.
\item The notes indicate replacing the Bernoulli/Beta update by a multinomial/Dirichlet update for four response categories.
\end{itemize}
}
